%%%%%%%%%%%%%%%%%%%%%%%%%%%%%%%%%%%%%%%%%%%%%%%%%%%%%%%%%%%%%%%%%%%%%%%%%%%
% Kapitel 0: Zum Modul -- Modulbeschreibung, Didaktik, Semesterplan,
% Leistungsnachweise, Gebrauch von Skript, Folien und Handouts.
% Quelle: WuZ-Buch, Kap. 15 (Curriculum, Modulhandbucheintrag),
%         Kap. 6 (Goldener Ablauf, Nocebo der Lehre), Anhang E (Grenzarbeit)
%%%%%%%%%%%%%%%%%%%%%%%%%%%%%%%%%%%%%%%%%%%%%%%%%%%%%%%%%%%%%%%%%%%%%%%%%%%
\nsl{Dieses Skript begleitet das Modul \emph{Wissenschaftliches Arbeiten} in
der Fassung \emph{Wissenschaft und Zauberkunst} (WuZ). Es setzt das
Curriculum aus Kapitel~15 des Buches \emph{Wissenschaft und Zauberkunst --
Wie wir uns ein besseres Bild von der Welt machen} in 13 Lehreinheiten um,
eine pro Semesterwoche. Jede Einheit beginnt mit einem Staunmoment, einer
kurzen Zaubervorführung oder einem Selbstversuch, dessen Wirkung das
Methodenthema der Woche vorwegnimmt. Die Hypothese, die dem Ansatz zugrunde
liegt, lautet: Staunen führt zu besserem Lernverhalten, objektiverem Denken
und weniger Verzerrung in Aussagen und Entscheidungen.
\newpage}

%==============================================================================
\section{Gebrauch von Skript, Folien und Handouts}
%==============================================================================
\nsl{Skript, Folien und Handouts entstehen aus denselben Quelldateien. Das
\textbf{Skript} (dieses Dokument) ist die vollständige Fassung für Lehrende:
Es enthält Fließtext, Folieninhalte, die Vorbereitung jedes Kunststücks, den
Ablauf der Doppelstunde, Hinweise zur Durchführung und die Erwartungshorizonte
der Übungen. Die \textbf{Folien} enthalten nur, was in der Stunde gezeigt
wird. Die \textbf{Handouts und Übungsblätter} sind eine eigene Druckvorlage
für die Studierenden, ohne Regie- und Lösungshinweise.

\begin{center}
\small
\begin{tabular}{@{}p{0.5\textwidth}ccc@{}}
\toprule
\textbf{Element} & \textbf{Skript} & \textbf{Folien} & \textbf{Handouts} \\
\midrule
Fließtext, Theorie & ja & -- & -- \\
Folieninhalte (Stichpunkte, Grafiken) & ja & ja & -- \\
\textcolor{rwuviolet}{\textbf{Staunmoment}} (violette Box) & ja & ja & -- \\
\textbf{Vorbereitung vor der Stunde} (graue Box) & ja & -- & -- \\
\textcolor{rwucyandark}{\textbf{Ablauf der Doppelstunde}} (cyanfarbene Box) & ja & -- & -- \\
Hinweis für Lehrende, Erwartungshorizont & ja & -- & -- \\
Handout (eine Seite je Einheit) & ja & -- & ja \\
Übungsblatt & ja & -- & ja \\
\bottomrule
\end{tabular}
\end{center}

Die Vorbereitungsboxen sind als Checklisten angelegt: was eine Woche
vorher zu besorgen und zu proben ist, was zu kopieren ist, was am Tag selbst
im Raum vorbereitet werden muss und was zu tun ist, wenn ein Effekt nicht
gelingt. Eine Gesamtliste aller Requisiten steht in Anhang~A.}

\begin{lehrhinweis}[Zur Geheimhaltung]
Dieses Skript ist öffentlich. Methoden kommerzieller und klassischer
Kunststücke (Switchboard, Chop Cup, Seilroutine, Becherroutinen) sind darum
nicht beschrieben; für sie gilt die Anleitung des Herstellers bzw. die
Standardliteratur. Beschrieben ist alles, was das Begleitbuch selbst
offenlegt: das Pendel, die Charpentier-Illusion, das Kruskal-Prinzip, die
Monty-Hall-Mathematik und der Vorhersage-Trick zum p-Hacking. Gegenüber den
Studierenden gilt die Haltung aus LE~13: Methoden werden nicht verraten,
die Richtung, in der sie liegen, darf angedeutet werden.
\end{lehrhinweis}

%==============================================================================
\section{Das Modul}
%==============================================================================
\begin{description}
\item[{\rot\bf Wissenschaftliches Arbeiten (WuZ)}]~\\[-\lblskip]
\bi
  \ite Grundstudium, Pflichtmodul, studiengangsübergreifend (Bachelor)
  \ite 5 ECTS, 150\,h Arbeitsaufwand (ca.\ 60\,h Präsenz, 90\,h
    Selbststudium und Projektarbeit)
  \ite seminaristische Vorlesung mit Demonstrationen, Kurzexperimenten,
    Einzel- und Gruppenübungen
  \ite semesterbegleitend: Tagebuch des Staunens und
    eigene Forschungsarbeit
  \ite Prüfung: Portfolio, benotet
\ei
\os{\newpage}

\item[{\rot\bf Was Sie am Ende können}]~\\[-\lblskip]
\bi
  \ite eine beantwortbare Forschungsfrage entwickeln und daraus prüfbare
    Hypothesen ableiten
  \ite Literatur systematisch recherchieren und Quellen kritisch bewerten
  \ite ein Forschungsdesign wählen, das ein echtes Widerlegungsrisiko eingeht
  \ite quantitative und qualitative Daten auswerten und Befunde von
    Deutungen trennen
  \ite Ergebnisse schreiben und präsentieren und dabei Unsicherheit offen
    benennen
  \ite die Regeln guter wissenschaftlicher Praxis auf die eigene Arbeit
    anwenden
\ei
\os{\newpage}
\end{description}

\nsl{Die Qualifikationsziele folgen den Kompetenzfeldern des
Qualifikationsrahmens für Deutsche Hochschulabschlüsse \parencite{hqr2017}:
\emph{Wissen und Verstehen} (Grundbegriffe und Ablauf wissenschaftlichen
Arbeitens, wissenschaftstheoretische Positionen, Grundlogik des
Signifikanztests und der bayesianischen Aktualisierung), \emph{Einsatz,
Anwendung und Erzeugung von Wissen} (eigene Forschungsfrage, Hypothesen,
Design und eine kleine kontrollierte Untersuchung), \emph{Nutzung und
Transfer} (Belastbarkeit von Quellen und Statistiken beurteilen),
\emph{Kommunikation und Kooperation} (Präsentation, Peer-Review,
Codier-Vergleich) sowie \emph{wissenschaftliches Selbstverständnis und
Professionalität} (reflexiv-selbstkritische Haltung, Normen
wissenschaftlicher Integrität). Der vollständige Modulhandbucheintrag steht
im Begleitbuch, Kapitel~15.}

%==============================================================================
\section{Der didaktische Ansatz}
%==============================================================================
\nsl{Jede Lehreinheit folgt dem Dreiklang \emph{Staunen erzeugen --
Erklärung suchen -- Methode verstehen}. Der Staunmoment ist der epistemische
Einstieg: Er erzeugt einen Prediction Error, also eine Abweichung zwischen
Erwartung und Erleben, die das Gehirn auflösen will \parencite{clark2016}.
Die Staun-Interventionen sind so gewählt, dass sie das Methodenthema
inhaltlich vorwegnehmen. Das Pendel in LE~1 zeigt, dass Wahrnehmung
Konstruktion ist; die Becher in LE~8 zeigen, dass das Sichtbare und der
Prozess dahinter auseinanderfallen können.}
\begin{description}
\item[{\rot\bf Jede Woche dasselbe Muster}]~\\[-\lblskip]
\bi
  \ite \textbf{Staunen}: ein Kunststück oder Selbstversuch, ohne Vorrede
  \ite \textbf{Erklärung}: Hypothesen sammeln, noch ohne Bewertung
  \ite \textbf{Methode}: das Werkzeug der Woche, mit dem man zwischen den
    Erklärungen entscheidet
  \ite \textbf{Anwendung}: Übung am eigenen Semesterprojekt
  \ite \textbf{Transfer}: Tagebuch und Aufgabe bis zur nächsten Woche
\ei
\os{\newpage}
\end{description}

\nsl{Für die Vorführung selbst gilt der \emph{Goldene Ablauf} aus dem
Begleitbuch (Kapitel~6): thematischer Einstieg, eine kleine Fertigkeit als
Glaubwürdigkeitssignal, der Effekt kurz und knapp, eine Aha-Pause von etwa
zehn Sekunden, dann der Transfer zum Thema. Zu vermeiden ist das
\enquote{Nocebo der Lehre}: der Satz \enquote{Das war jetzt nur ein Trick},
der das Staunen entwertet, bevor es wirken konnte. Das Prinzip hinter dem
Effekt wird erklärt, das Gefühl des Nicht-Wissens darf aber einen Moment
stehen bleiben. Glaubwürdigkeit entsteht dabei aus Handlungen, nicht aus
Behauptungen \parencite{henrich2009}; wer ein Kunststück vorführt, sollte es
deshalb so oft geübt haben, dass es sicher sitzt.}

\begin{lehrhinweis}[Goldener Ablauf einer Vorführung]
\begin{enumerate}
  \item Thematischer Einstieg (ein Satz, der das Thema der Woche anreißt)
  \item eine kleine Fertigkeit oder ein besonderes Requisit
  \item der Effekt, kurz und knapp
  \item Aha-Pause: das Staunen etwa zehn Sekunden wirken lassen
  \item Transfer zum Methodenthema
\end{enumerate}
\end{lehrhinweis}

%==============================================================================
\section{Semesterplan}
%==============================================================================
\begin{description}
\item[{\rot\bf 13 Lehreinheiten}]~\\[-\lblskip]
\begin{center}
\nsl{\small}\os{\scriptsize\renewcommand{\arraystretch}{0.95}}
\rowcolors{2}{rwucyan40!20}{white}
\begin{tabularx}{\linewidth}{@{}>{\bfseries\color{rwuviolet}}p{6mm}>{\raggedright\arraybackslash}p{0.3\textwidth}X>{\raggedright\arraybackslash}p{0.22\textwidth}@{}}
\toprule
\rowcolor{rwuviolet}\color{white}LE & \color{white}\textbf{Thema} & \color{white}\textbf{Staunmoment} & \color{white}\textbf{Kernbegriff} \\
\midrule
1  & Staunen als Erkenntnishaltung     & Ideomotor-Pendel                      & Prediction Error \\
2  & Die Forschungsfrage               & Das geheime Objekt                    & Operationalisierung \\
3  & Literaturrecherche                & Kruskal Count                         & Suchstrategie \\
4  & Quellenkritik                     & Apophenie / Das letzte Abendmahl      & Apophenie / Bias \\
5  & Hypothesen und Falsifizierbarkeit & Switchboard / Black Box               & Falsifikation \\
6  & Forschungsdesign                  & Kontrollierte Seilroutine             & Experimentelle Kontrolle \\
7  & Kognitive Verzerrungen            & Charpentier-Illusion                  & Confirmation Bias \\
8  & Quantitative Methoden             & Chop Cup / Monty Hall                 & Bedingte Wahrscheinlichkeit \\
9  & Qualitative Methoden              & Das Patter                            & Intersubjektivität \\
10 & Datenanalyse und Interpretation   & Die unveränderliche Summe             & Determinismus \\
11 & Wissenschaftliches Schreiben      & Patter-Umbau in Fachsprache           & Präzision und Struktur \\
12 & Präsentation und Glaubwürdigkeit  & CRED-Strategien                       & Credibility Cues \\
13 & Wissenschaftsethik und Grenzarbeit & Boundary Work: Was verrät man?       & Forschungsethik \\
\bottomrule
\end{tabularx}
\end{center}
\os{\newpage}
\end{description}

%==============================================================================
\section{Semesterprojekt und Leistungsnachweise}
%==============================================================================
\nsl{Parallel zu den Lehreinheiten bearbeitet jede Person ein eigenes kleines
Forschungsvorhaben. Die Teilleistungen bauen aufeinander auf: Wer
Rechercheprotokoll, Design-Skizze und Präsentation sorgfältig erarbeitet hat,
hat einen großen Teil der Seminararbeit bereits geschrieben. Die folgende
Grafik zeigt, in welcher Einheit welcher Baustein entsteht.}
\begin{description}
\item[{\rot\bf Der rote Faden: Ihr Projekt}]~\\[-\lblskip]
\begin{center}
\resizebox{\linewidth}{!}{%
\begin{tikzpicture}[font=\sffamily\small,>=Stealth,
  st/.style={draw=rwuviolet,thick,fill=rwuviolet!8,rounded corners=2pt,
    minimum width=24mm,minimum height=12mm,align=center},
  ab/.style={draw=rwucyan,thick,fill=rwucyan40!40,rounded corners=2pt,
    minimum width=24mm,minimum height=12mm,align=center}]
  \node[st] (a) at (0,0)    {LE 2\\Forschungsfrage};
  \node[ab] (b) at (3,0)    {LE 3\\Recherche-\\protokoll};
  \node[st] (c) at (6,0)    {LE 5\\Hypothesen};
  \node[ab] (d) at (9,0)    {LE 6\\Design-Skizze};
  \node[st] (e) at (12,0)   {LE 8--10\\Daten und\\Auswertung};
  \node[ab] (f) at (15,0)   {LE 12\\Präsentation};
  \node[ab,draw=rwuviolet,fill=rwuviolet!25] (g) at (18,0) {Semester-\\ende\\Seminararbeit};
  \foreach \x/\y in {a/b,b/c,c/d,d/e,e/f,f/g} \draw[->,thick] (\x) -- (\y);
  \node[font=\sffamily\footnotesize,text=rwugray] at (9,-1.4) {cyan: Teilleistung im Portfolio \quad violett: Arbeitsschritt im Projekt};
\end{tikzpicture}}
\end{center}
\os{\newpage}

\item[{\rot\bf Portfolio}]~\\[-\lblskip]
\begin{center}
\small
\rowcolors{2}{rwucyan40!20}{white}
\begin{tabularx}{\linewidth}{@{}>{\bfseries}p{0.27\textwidth}p{0.15\textwidth}Xr@{}}
\toprule
\rowcolor{rwuviolet}\color{white}Leistung & \color{white}\textbf{Fällig} & \color{white}\textbf{Beschreibung} & \color{white}\textbf{Anteil} \\
\midrule
Tagebuch des Staunens & Semesterende & Reflexionsheft, 13 Einträge & 20\,\% \\
Rechercheprotokoll & LE~3 & dokumentierte Recherche mit Quellenliste & 10\,\% \\
Forschungsdesign-Skizze & LE~6 & 1--2 Seiten, begründetes Studiendesign & 10\,\% \\
5-Minuten-Präsentation & LE~12 & Zwischenbericht zum Semesterprojekt & 15\,\% \\
Seminararbeit & Semesterende & 12--15 Seiten, vollständiger Forschungsbericht & 45\,\% \\
\bottomrule
\end{tabularx}
\end{center}
\os{\newpage}

\item[{\rot\bf Das Tagebuch des Staunens}]~\\[-\lblskip]
\bi
  \ite ein Heft, das Sie ab LE~1 führen: jede Woche ein Eintrag
  \ite (1) Was habe ich erwartet? (2) Was ist passiert? (3) Welcher Begriff
    aus der Vorlesung passt dazu?
  \ite Ziel: Staunen bemerken und mit einem wissenschaftlichen Begriff
    verknüpfen
  \ite letzter Eintrag in LE~13: Was hat mich in diesem Semester am meisten
    überrascht?
\ei
\os{\newpage}
\end{description}

\begin{lehrhinweis}[Vor Semesterbeginn]
\begin{checkliste}
  \item Anhang~A durchgehen und Requisiten für das ganze Semester bestellen;
    Lieferzeiten für kommerzielle Requisiten (Switchboard, Seilset, Chop Cup)
    einplanen.
  \item Kunststücke mindestens zwei Wochen vor dem jeweiligen Termin proben;
    das Switchboard braucht die längste Einübung.
  \item \texttt{handouts.pdf} pro Woche kopieren; Anzahl nach Anmeldestand
    plus Reserve.
  \item A5-Hefte für das Tagebuch des Staunens bereitlegen oder in der
    Ankündigung um ein eigenes Heft bitten.
  \item Raum prüfen: Tafel oder Whiteboard, Beamer, Tische für Gruppenarbeit,
    Platz zum Herumreichen von Gegenständen.
\end{checkliste}
\end{lehrhinweis}
